% Options for packages loaded elsewhere
\PassOptionsToPackage{unicode}{hyperref}
\PassOptionsToPackage{hyphens}{url}
%
\documentclass[
]{article}
\usepackage{amsmath,amssymb}
\usepackage{iftex}
\ifPDFTeX
  \usepackage[T1]{fontenc}
  \usepackage[utf8]{inputenc}
  \usepackage{textcomp} % provide euro and other symbols
\else % if luatex or xetex
  \usepackage{unicode-math} % this also loads fontspec
  \defaultfontfeatures{Scale=MatchLowercase}
  \defaultfontfeatures[\rmfamily]{Ligatures=TeX,Scale=1}
\fi
\usepackage{lmodern}
\ifPDFTeX\else
  % xetex/luatex font selection
\fi
% Use upquote if available, for straight quotes in verbatim environments
\IfFileExists{upquote.sty}{\usepackage{upquote}}{}
\IfFileExists{microtype.sty}{% use microtype if available
  \usepackage[]{microtype}
  \UseMicrotypeSet[protrusion]{basicmath} % disable protrusion for tt fonts
}{}
\makeatletter
\@ifundefined{KOMAClassName}{% if non-KOMA class
  \IfFileExists{parskip.sty}{%
    \usepackage{parskip}
  }{% else
    \setlength{\parindent}{0pt}
    \setlength{\parskip}{6pt plus 2pt minus 1pt}}
}{% if KOMA class
  \KOMAoptions{parskip=half}}
\makeatother
\usepackage{xcolor}
\usepackage{longtable,booktabs,array}
\usepackage{calc} % for calculating minipage widths
% Correct order of tables after \paragraph or \subparagraph
\usepackage{etoolbox}
\makeatletter
\patchcmd\longtable{\par}{\if@noskipsec\mbox{}\fi\par}{}{}
\makeatother
% Allow footnotes in longtable head/foot
\IfFileExists{footnotehyper.sty}{\usepackage{footnotehyper}}{\usepackage{footnote}}
\makesavenoteenv{longtable}
\setlength{\emergencystretch}{3em} % prevent overfull lines
\providecommand{\tightlist}{%
  \setlength{\itemsep}{0pt}\setlength{\parskip}{0pt}}
\setcounter{secnumdepth}{-\maxdimen} % remove section numbering
\ifLuaTeX
  \usepackage{selnolig}  % disable illegal ligatures
\fi
\IfFileExists{bookmark.sty}{\usepackage{bookmark}}{\usepackage{hyperref}}
\IfFileExists{xurl.sty}{\usepackage{xurl}}{} % add URL line breaks if available
\urlstyle{same}
\hypersetup{
  hidelinks,
  pdfcreator={LaTeX via pandoc}}

\author{}
\date{}

\begin{document}

\hypertarget{ayudantuxeda-8-repaso-i2-pauta-completa}{%
\section{Ayudantía 8 --- Repaso I2: Pauta
Completa}\label{ayudantuxeda-8-repaso-i2-pauta-completa}}

\textbf{Pontificia Universidad Católica de Chile} \textbf{Departamento
Ingeniería Industrial y de Sistemas} \textbf{ICS3213 -- Gestión de
Operaciones} \emph{Profesores: Alejandro Mac Cawley --- Rodrigo Carrasco
(Secciones 1, 2 y 3)} \emph{Primer Semestre 2026} \emph{Ayudante: Juan
Pablo García --- jgarca@uc.cl}

\begin{center}\rule{0.5\linewidth}{0.5pt}\end{center}

\hypertarget{repaso-teuxf3rico-y-consejos-por-muxf3dulo}{%
\subsection{Repaso Teórico y Consejos por
Módulo}\label{repaso-teuxf3rico-y-consejos-por-muxf3dulo}}

\hypertarget{muxf3dulo-1-planificaciuxf3n-agregada}{%
\subsubsection{Módulo 1: Planificación
Agregada}\label{muxf3dulo-1-planificaciuxf3n-agregada}}

La planificación agregada traduce la estrategia de una empresa en
acciones concretas, actuando como puente entre la estrategia competitiva
y las operaciones diarias.

\textbf{Nivel de Decisión:} Táctico con impacto operativo (mediano
plazo), entre la planeación estratégica (largo plazo) y la programación
diaria (corto plazo).

\textbf{Concepto de ``Agregación'':} - Los pronósticos por producto
individual (SKU) tienen alto nivel de error. - Al consolidar la demanda
de múltiples productos, los errores individuales tienden a compensarse
(ley de los grandes números). - Se usan ``unidades equivalentes'' que
agrupan productos similares. - Objetivo: igualar demanda prevista con
capacidad productiva, tomando decisiones sobre producción,
subcontratación, inventario y fuerza laboral.

\textbf{Estrategias para abordar la demanda:}

\begin{longtable}[]{@{}
  >{\raggedright\arraybackslash}p{(\columnwidth - 6\tabcolsep) * \real{0.2500}}
  >{\raggedright\arraybackslash}p{(\columnwidth - 6\tabcolsep) * \real{0.2500}}
  >{\raggedright\arraybackslash}p{(\columnwidth - 6\tabcolsep) * \real{0.2500}}
  >{\raggedright\arraybackslash}p{(\columnwidth - 6\tabcolsep) * \real{0.2500}}@{}}
\toprule\noalign{}
\begin{minipage}[b]{\linewidth}\raggedright
Estrategia
\end{minipage} & \begin{minipage}[b]{\linewidth}\raggedright
Descripción
\end{minipage} & \begin{minipage}[b]{\linewidth}\raggedright
Ventaja
\end{minipage} & \begin{minipage}[b]{\linewidth}\raggedright
Desventaja
\end{minipage} \\
\midrule\noalign{}
\endhead
\bottomrule\noalign{}
\endlastfoot
\textbf{Persecución (Chase)} & Ajusta capacidad para seguir la demanda
período a período (contratar/despedir/horas extra) & Inventario mínimo &
Costos altos de contratación/despido, impacto en moral \\
\textbf{Nivelación (Level)} & Tasa de producción constante; el desajuste
se absorbe con inventario o faltantes & Fuerza laboral estable, sin
costos de rotación & Costos de inventario o pérdida de ventas \\
\textbf{Mixta} & Combinación de ambas según temporada, costo y
restricciones laborales & Equilibrio costo-flexibilidad & Mayor
complejidad de planificación \\
\end{longtable}

\begin{center}\rule{0.5\linewidth}{0.5pt}\end{center}

\hypertarget{muxf3dulo-2-mrp-planificaciuxf3n-de-requerimientos-de-materiales}{%
\subsubsection{Módulo 2: MRP (Planificación de Requerimientos de
Materiales)}\label{muxf3dulo-2-mrp-planificaciuxf3n-de-requerimientos-de-materiales}}

\textbf{Jerarquía de planificación:}
\[\text{Planificación Agregada} \to \text{MPS} \to \text{MRP} \to \text{Programación diaria}\]

\begin{itemize}
\tightlist
\item
  \textbf{MPS:} Traduce el plan agregado en productos específicos,
  cantidades y períodos. Responde: ¿cuántas unidades del producto X en
  la semana Y?
\item
  \textbf{MRP:} Explota el MPS hacia atrás en el tiempo usando el BOM y
  los Lead Times. Responde: ¿qué materiales, en qué cantidad y cuándo
  lanzar órdenes?
\end{itemize}

\textbf{Métodos de Lotificación:}

\begin{longtable}[]{@{}
  >{\raggedright\arraybackslash}p{(\columnwidth - 6\tabcolsep) * \real{0.2500}}
  >{\raggedright\arraybackslash}p{(\columnwidth - 6\tabcolsep) * \real{0.2500}}
  >{\raggedright\arraybackslash}p{(\columnwidth - 6\tabcolsep) * \real{0.2500}}
  >{\raggedright\arraybackslash}p{(\columnwidth - 6\tabcolsep) * \real{0.2500}}@{}}
\toprule\noalign{}
\begin{minipage}[b]{\linewidth}\raggedright
Método
\end{minipage} & \begin{minipage}[b]{\linewidth}\raggedright
Lógica
\end{minipage} & \begin{minipage}[b]{\linewidth}\raggedright
Ventaja
\end{minipage} & \begin{minipage}[b]{\linewidth}\raggedright
Desventaja
\end{minipage} \\
\midrule\noalign{}
\endhead
\bottomrule\noalign{}
\endlastfoot
\textbf{Lote a Lote (L×L)} & Ordena exactamente lo necesario cada
período & Cero inventario & Muchos setups si la demanda es frecuente \\
\textbf{Cantidad Fija} & Siempre la misma cantidad predefinida & Simple
& Puede generar inventario innecesario \\
\textbf{Período Fijo} & Ordena cada N períodos & Reduce frecuencia de
órdenes & Inventario variable e impredecible \\
\textbf{EOQ} & Cantidad óptima con demanda promedio constante & Óptimo
con demanda estable & Asume demanda constante (rara vez real en MRP) \\
\textbf{Silver-Meal} & Minimiza costo promedio por período
heurísticamente & Se adapta a demanda variable, simple & No garantiza el
óptimo global \\
\textbf{Wagner-Whitin} & Programación dinámica exacta & Óptimo
matemático & Computacionalmente costoso, menos intuitivo \\
\end{longtable}

\textbf{Consejo clave para Silver-Meal:} Calcular \(C(1), C(2), \ldots\)
hasta que \(C(k+1) > C(k)\), donde:
\[C(k) = \frac{S + H \sum_{j=1}^{k-1} j \cdot D_{t+j}}{k}\] Parar en
\(k\) y lanzar ese lote. Iniciar nueva iteración desde el período
siguiente con demanda.

\begin{center}\rule{0.5\linewidth}{0.5pt}\end{center}

\hypertarget{muxf3dulo-3-pert-programaciuxf3n-de-proyectos}{%
\subsubsection{Módulo 3: PERT (Programación de
Proyectos)}\label{muxf3dulo-3-pert-programaciuxf3n-de-proyectos}}

\textbf{Distribución Beta para duración de actividad \((a, m, b)\):}
\[t_e = \frac{a + 4m + b}{6}, \qquad \sigma = \frac{b - a}{6}, \qquad \sigma^2 = \left(\frac{b-a}{6}\right)^2\]

\textbf{Tiempos de calendario:} - \textbf{ES} (Earliest Start): máximo
de los EF de todos los predecesores. Para actividades sin predecesores:
ES = 0. - \textbf{EF} (Earliest Finish): \(EF = ES + t_e\). -
\textbf{LF} (Latest Finish): mínimo de los LS de los sucesores. Para la
última actividad: \(LF = T\). - \textbf{LS} (Latest Start):
\(LS = LF - t_e\). - \textbf{Holgura:} \(H = LS - ES = LF - EF\). -
\textbf{Ruta Crítica:} actividades con holgura = 0. La duración del
proyecto es \(T = \max(EF)\).

\textbf{Estadísticas de la ruta crítica:}
\[\mu_p = \sum_{i \in RC} t_{e,i}, \qquad \sigma_p^2 = \sum_{i \in RC} \sigma_i^2, \qquad \sigma_p = \sqrt{\sigma_p^2}\]

\textbf{Crashing:} Reducción de duración a costo mínimo. La pendiente de
crashing por actividad es:
\[\text{Pendiente}_i = \frac{CC_i - CN_i}{TN_i - TC_i} \quad [\$/\text{semana}]\]
Menor pendiente = más barato. Evaluar por valor esperado neto:
\(VE = \text{bono} \cdot P(\text{terminar antes}) - \text{costo aceleración}\).

\begin{center}\rule{0.5\linewidth}{0.5pt}\end{center}

\hypertarget{muxf3dulo-4-variabilidad-teoruxeda-de-colas}{%
\subsubsection{Módulo 4: Variabilidad (Teoría de
Colas)}\label{muxf3dulo-4-variabilidad-teoruxeda-de-colas}}

\textbf{Ley de Little:}
\(L = \lambda \cdot W \implies WIP = TH \cdot CT\)

\textbf{Coeficiente de variación:} \(c_x = \sigma_x / \mu_x\)

\textbf{Disponibilidad ante fallas:} \[A = \frac{MTBF}{MTBF + MTTR}\]
\textgreater{} Paradas largas e infrecuentes aumentan \textbf{mucho más}
la variabilidad efectiva que paradas cortas y frecuentes.

\textbf{Impacto de Setups} en tiempo efectivo y varianza:
\[t_e = t_o + \frac{t_s}{N_s}, \qquad \sigma_e^2 = \sigma_o^2 + \frac{\sigma_s^2}{N_s} + \frac{N_s - 1}{N_s^2} t_s^2\]

\textbf{Modelo G/G/1 --- Ecuación de Kingman (VUT):}
\[CT_q = \left(\frac{c_a^2 + c_e^2}{2}\right) \cdot \frac{\rho}{1-\rho} \cdot t_e, \qquad \rho = \frac{\lambda}{\mu}\]

\textbf{Propagación de variabilidad en serie:}
\[c_s^2 \approx \rho^2 \cdot c_e^2 + (1 - \rho^2) \cdot c_a^2\]

\textbf{Estrategias prácticas:} - La utilización \(\rho\) es el factor
crítico: cuando \(\rho \to 1\), el tiempo en cola explota
exponencialmente. - \textbf{Pooling (unifila):} consolida recursos y
promedia la variabilidad global --- siempre reduce la espera. -
\textbf{Psicología de la espera:} el cliente percibe el tiempo ocioso y
las esperas injustas/inciertas como más largas. Distraer, explicar
retrasos, ocultar empleados inactivos.

\begin{center}\rule{0.5\linewidth}{0.5pt}\end{center}

\hypertarget{problema-1-planificaciuxf3n-agregada-la-miga-dorada}{%
\subsection{Problema 1: Planificación Agregada --- La Miga
Dorada}\label{problema-1-planificaciuxf3n-agregada-la-miga-dorada}}

\hypertarget{enunciado}{%
\subsubsection{Enunciado}\label{enunciado}}

La panadería ``La Miga Dorada'' elabora panes artesanales. Cada pan
requiere 1 hora-hombre (HH) de producción. Demanda mensual del primer
trimestre:

\begin{longtable}[]{@{}lccc@{}}
\toprule\noalign{}
Mes & Enero & Febrero & Marzo \\
\midrule\noalign{}
\endhead
\bottomrule\noalign{}
\endlastfoot
\textbf{Demanda (panes)} & 1 200 & 2 000 & 1 000 \\
\end{longtable}

\textbf{Datos:} - Dotación inicial: 10 panaderos. - Días hábiles por
mes: 20 días × 8 horas = \textbf{160 HH/mes por panadero}. - Costo HH
normal: \(\$50\)/hora. - Costo HH extra: \(\$80\)/hora. - Costo
inventario: \(\$10\)/pan al mes. - Costo faltante: \(\$15\)/pan al mes.
- Contratar: \(\$1\,500\)/panadero. Despedir: \(\$2\,500\)/panadero.

\textbf{Se pide comparar:} 1. \textbf{Plan Chase:} ajustar panaderos mes
a mes para producir exactamente la demanda. 2. \textbf{Plan Nivel con
Inventario y Faltante:} mantener 10 panaderos; producir máximo normal;
permitir inventario o faltantes. 3. \textbf{Plan Nivel con Horas Extra:}
mantener 10 panaderos; cubrir exceso con horas extra; sin faltantes ni
inventario.

\hypertarget{soluciuxf3n-oficial}{%
\subsubsection{Solución Oficial}\label{soluciuxf3n-oficial}}

\textbf{Datos previos:} - Capacidad normal mensual (10 panaderos):
\(10 \times 160 = 1\,600\) panes/mes.

\begin{center}\rule{0.5\linewidth}{0.5pt}\end{center}

\hypertarget{plan-1-chase}{%
\paragraph{Plan 1 --- Chase}\label{plan-1-chase}}

Panaderos requeridos cada mes = demanda / 160, \textbf{redondeando al
entero más cercano}:

\begin{longtable}[]{@{}
  >{\raggedright\arraybackslash}p{(\columnwidth - 8\tabcolsep) * \real{0.1739}}
  >{\centering\arraybackslash}p{(\columnwidth - 8\tabcolsep) * \real{0.2174}}
  >{\centering\arraybackslash}p{(\columnwidth - 8\tabcolsep) * \real{0.2174}}
  >{\raggedright\arraybackslash}p{(\columnwidth - 8\tabcolsep) * \real{0.1739}}
  >{\centering\arraybackslash}p{(\columnwidth - 8\tabcolsep) * \real{0.2174}}@{}}
\toprule\noalign{}
\begin{minipage}[b]{\linewidth}\raggedright
Mes
\end{minipage} & \begin{minipage}[b]{\linewidth}\centering
Demanda
\end{minipage} & \begin{minipage}[b]{\linewidth}\centering
Panaderos req.
\end{minipage} & \begin{minipage}[b]{\linewidth}\raggedright
Transición
\end{minipage} & \begin{minipage}[b]{\linewidth}\centering
Costo cont./desp.
\end{minipage} \\
\midrule\noalign{}
\endhead
\bottomrule\noalign{}
\endlastfoot
Enero & 1 200 & 8 & 10 → 8: \textbf{despido de 2} &
\(2 \times 2\,500 = \$5\,000\) \\
Febrero & 2 000 & 13 & 8 → 13: \textbf{contratación de 5} &
\(5 \times 1\,500 = \$7\,500\) \\
Marzo & 1 000 & 7 & 13 → 7: \textbf{despido de 6} &
\(6 \times 2\,500 = \$15\,000\) \\
\textbf{Total} & & & & \textbf{\(\$27\,500\)} \\
\end{longtable}

\begin{itemize}
\tightlist
\item
  Costo MO normal:
  \((1\,200 + 2\,000 + 1\,000) \times \$50 = 4\,200 \times 50 = \$210\,000\).
\item
  Inventario / Faltante: \(\$0\) (produce exactamente la demanda).
\end{itemize}

\[\boxed{\text{Costo Total Plan Chase} = 210\,000 + 27\,500 = \$237\,500}\]

\begin{quote}
\textbf{Nota de cálculo:} La pauta redondea a enteros al calcular los
panaderos requeridos (\(1200/160 = 7{,}5 \to 8\), etc.). Si se trabaja
con fracciones el costo de cont./desp. cambia ligeramente pero la lógica
es idéntica.
\end{quote}

\begin{center}\rule{0.5\linewidth}{0.5pt}\end{center}

\hypertarget{plan-2-nivel-con-inventario-y-faltante}{%
\paragraph{Plan 2 --- Nivel con Inventario y
Faltante}\label{plan-2-nivel-con-inventario-y-faltante}}

10 panaderos fijos → producción = 1 600 panes/mes:

\begin{longtable}[]{@{}lccccc@{}}
\toprule\noalign{}
Mes & Producción & Demanda & Inv. final & Faltante & Costo \\
\midrule\noalign{}
\endhead
\bottomrule\noalign{}
\endlastfoot
Enero & 1 600 & 1 200 & \textbf{400} & 0 &
\(400 \times 10 = \$4\,000\) \\
Febrero & 1 600 & 2 000 & \textbf{0} & 0 & \(\$0\) \\
Marzo & 1 600 & 1 000 & \textbf{600} & 0 &
\(600 \times 10 = \$6\,000\) \\
\end{longtable}

\begin{itemize}
\tightlist
\item
  Costo MO normal: \(3 \times 1\,600 \times \$50 = \$240\,000\).
\item
  Costo inventario total: \(\$10\,000\).
\end{itemize}

\[\boxed{\text{Costo Total Plan Nivel + Inv./Falt.} = 240\,000 + 10\,000 = \$250\,000}\]

\begin{center}\rule{0.5\linewidth}{0.5pt}\end{center}

\hypertarget{plan-3-nivel-con-horas-extra}{%
\paragraph{Plan 3 --- Nivel con Horas
Extra}\label{plan-3-nivel-con-horas-extra}}

10 panaderos fijos → producción normal máxima = 1 600 panes/mes; el
exceso se cubre con HH extra:

\begin{longtable}[]{@{}lccccc@{}}
\toprule\noalign{}
Mes & Demanda & Prod. normal & Exceso & HH extra & Costo HE \\
\midrule\noalign{}
\endhead
\bottomrule\noalign{}
\endlastfoot
Enero & 1 200 & 1 200 & 0 & 0 & \(\$0\) \\
Febrero & 2 000 & 1 600 & 400 & 400 & \(400 \times 80 = \$32\,000\) \\
Marzo & 1 000 & 1 000 & 0 & 0 & \(\$0\) \\
\end{longtable}

\begin{itemize}
\tightlist
\item
  Costo MO normal: \(\$240\,000\).
\item
  Costo HH extra: \(\$32\,000\).
\end{itemize}

\[\boxed{\text{Costo Total Plan Nivel + HE} = 240\,000 + 32\,000 = \$272\,000}\]

\begin{center}\rule{0.5\linewidth}{0.5pt}\end{center}

\hypertarget{conclusiuxf3n}{%
\paragraph{Conclusión}\label{conclusiuxf3n}}

\begin{longtable}[]{@{}lc@{}}
\toprule\noalign{}
Plan & Costo Total \\
\midrule\noalign{}
\endhead
\bottomrule\noalign{}
\endlastfoot
Chase & \(\$237\,500\) \\
Nivel con Inventario y Faltante & \(\$250\,000\) \\
Nivel con Horas Extra & \(\$272\,000\) \\
\end{longtable}

\textbf{El plan más económico es Plan Chase} (\(\$237\,500\)).

\begin{quote}
\textbf{Supuesto de redondeo:} la pauta redondea los panaderos
requeridos a enteros (\(1200/160 = 7{,}5 \to 8\)). Esto significa que en
enero se producen \(8 \times 160 = 1\,280\) panes, no 1,200 exactos ---
el plan Chase no es estrictamente puro. En la prueba, si el enunciado no
especifica, declara tu supuesto (enteros o fracciones) y sé consistente.
Ambos son aceptados si el cálculo es coherente.
\end{quote}

\begin{center}\rule{0.5\linewidth}{0.5pt}\end{center}

\hypertarget{problema-2-planificaciuxf3n-cp-y-mrp}{%
\subsection{Problema 2: Planificación CP y
MRP}\label{problema-2-planificaciuxf3n-cp-y-mrp}}

\hypertarget{enunciado-1}{%
\subsubsection{Enunciado}\label{enunciado-1}}

Una empresa fabrica un producto final ensamblado a partir de: - 1 unidad
de A → requiere 1 unidad de A1 y 2 de A2. - 2 unidades de B → cada B
requiere 1 unidad de A1 y 1 de B2. - 1 unidad de C.

\textbf{Lead Times:} Ensamblaje final = 1 semana. Todos los demás
componentes = 2 semanas.

Demanda del producto final (semanas 1--8):

\begin{longtable}[]{@{}lcccccccc@{}}
\toprule\noalign{}
Período & 1 & 2 & 3 & 4 & 5 & 6 & 7 & 8 \\
\midrule\noalign{}
\endhead
\bottomrule\noalign{}
\endlastfoot
\textbf{Demanda} & 0 & 0 & 0 & 10 & 50 & 40 & 60 & 50 \\
\end{longtable}

\textbf{Se pide:} 1. Árbol BOM. 2. Matriz MRP del producto final con
inventario inicial = 50 unidades, orden en tránsito = 20 unidades con
llegada en semana 2. Usar \textbf{Lote a Lote (L4L)}. 3. Demanda
consolidada de A1:

\begin{longtable}[]{@{}lcccccccc@{}}
\toprule\noalign{}
Período & 1 & 2 & 3 & 4 & 5 & 6 & 7 & 8 \\
\midrule\noalign{}
\endhead
\bottomrule\noalign{}
\endlastfoot
\textbf{Demanda A1} & 0 & 30 & 120 & 180 & 150 & 0 & 0 & 0 \\
\end{longtable}

Costo de setup \(S = \$120\), costo de inventario \(H = \$0{,}9\) por
unidad/semana. Usar \textbf{Silver-Meal}.

\hypertarget{soluciuxf3n-oficial-1}{%
\subsubsection{Solución Oficial}\label{soluciuxf3n-oficial-1}}

\hypertarget{i-uxe1rbol-bom}{%
\paragraph{i) Árbol BOM}\label{i-uxe1rbol-bom}}

\begin{verbatim}
Producto Final
├── A  (×1)
│   ├── A1 (×1)
│   └── A2 (×2)
├── B  (×2)
│   ├── A1 (×1)
│   └── B2 (×1)
└── C  (×1)
\end{verbatim}

\begin{quote}
\textbf{Nota:} A1 aparece en dos ramas: lo necesita A (×1) y B (×1 por
cada B, pero se usan 2 B por PF → 2 unidades de A1 provenientes de B).
La demanda consolidada de A1 =
\(1 \times \text{lanz. A} + 1 \times \text{lanz. B}\).
\end{quote}

\begin{center}\rule{0.5\linewidth}{0.5pt}\end{center}

\hypertarget{ii-matriz-mrp-producto-final-l4l-lt-1-semana}{%
\paragraph{ii) Matriz MRP --- Producto Final (L4L, LT = 1
semana)}\label{ii-matriz-mrp-producto-final-l4l-lt-1-semana}}

\begin{itemize}
\tightlist
\item
  Inventario inicial (período 0): \textbf{50 unidades}.
\item
  Orden en tránsito (SR): \textbf{20 unidades} en semana 2.
\end{itemize}

\begin{longtable}[]{@{}
  >{\raggedright\arraybackslash}p{(\columnwidth - 18\tabcolsep) * \real{0.0816}}
  >{\centering\arraybackslash}p{(\columnwidth - 18\tabcolsep) * \real{0.1020}}
  >{\centering\arraybackslash}p{(\columnwidth - 18\tabcolsep) * \real{0.1020}}
  >{\centering\arraybackslash}p{(\columnwidth - 18\tabcolsep) * \real{0.1020}}
  >{\centering\arraybackslash}p{(\columnwidth - 18\tabcolsep) * \real{0.1020}}
  >{\centering\arraybackslash}p{(\columnwidth - 18\tabcolsep) * \real{0.1020}}
  >{\centering\arraybackslash}p{(\columnwidth - 18\tabcolsep) * \real{0.1020}}
  >{\centering\arraybackslash}p{(\columnwidth - 18\tabcolsep) * \real{0.1020}}
  >{\centering\arraybackslash}p{(\columnwidth - 18\tabcolsep) * \real{0.1020}}
  >{\centering\arraybackslash}p{(\columnwidth - 18\tabcolsep) * \real{0.1020}}@{}}
\toprule\noalign{}
\begin{minipage}[b]{\linewidth}\raggedright
Período
\end{minipage} & \begin{minipage}[b]{\linewidth}\centering
0
\end{minipage} & \begin{minipage}[b]{\linewidth}\centering
1
\end{minipage} & \begin{minipage}[b]{\linewidth}\centering
2
\end{minipage} & \begin{minipage}[b]{\linewidth}\centering
3
\end{minipage} & \begin{minipage}[b]{\linewidth}\centering
4
\end{minipage} & \begin{minipage}[b]{\linewidth}\centering
5
\end{minipage} & \begin{minipage}[b]{\linewidth}\centering
6
\end{minipage} & \begin{minipage}[b]{\linewidth}\centering
7
\end{minipage} & \begin{minipage}[b]{\linewidth}\centering
8
\end{minipage} \\
\midrule\noalign{}
\endhead
\bottomrule\noalign{}
\endlastfoot
\textbf{Req. Bruto (GR)} & & 0 & 0 & 0 & 10 & 50 & 40 & 60 & 50 \\
\textbf{Recepciones Program. (SR)} & & 0 & 20 & 0 & 0 & 0 & 0 & 0 & 0 \\
\textbf{Inventario Final (OH)} & 50 & 50 & 70 & 70 & 60 & 10 & 0 & 0 &
0 \\
\textbf{Req. Neto (NR)} & & 0 & 0 & 0 & 0 & 0 & 30 & 60 & 50 \\
\textbf{Recepción Planeada (POR)} & & 0 & 0 & 0 & 0 & 0 & 30 & 60 &
50 \\
\textbf{Lanzamiento Planeado (PORelease)} & & 0 & 0 & 0 & 0 & 30 & 60 &
50 & --- \\
\end{longtable}

\begin{quote}
\textbf{Cálculo inventario:} \(OH_t = OH_{t-1} + SR_t + POR_t - GR_t\).
Ej.: \(OH_2 = 50 + 20 + 0 - 0 = 70\).
\end{quote}

\begin{center}\rule{0.5\linewidth}{0.5pt}\end{center}

\hypertarget{iii-silver-meal-para-a1}{%
\paragraph{iii) Silver-Meal para A1}\label{iii-silver-meal-para-a1}}

\(S = \$120\), \(H = \$0{,}9\)/unidad/semana. Demanda:
\([0, 30, 120, 180, 150, 0, 0, 0]\) (semanas 1--8).

\begin{quote}
\textbf{Regla de aplicación:} Silver-Meal arranca siempre en el
\textbf{primer período con demanda positiva}. El período 1 tiene demanda
= 0: no hay nada que ordenar, \(Q_1 = 0\) trivialmente. Empezamos el
algoritmo en el período 2.
\end{quote}

\textbf{Lote 1 --- arranca en período 2 (\(D_2 = 30\)):} - \(k=1\) (solo
sem. 2): \(C(1) = 120\). - \(k=2\) (sem. 2--3, \(D_3 = 120\)):
\(C(2) = (120 + 1 \cdot 0{,}9 \cdot 120)/2 = 228/2 = 114 < 120\) →
seguir. - \(k=3\) (sem. 2--4, \(D_4 = 180\)):
\(C(3) = (228 + 2 \cdot 0{,}9 \cdot 180)/3 = 552/3 = 184 > 114\) →
\textbf{PARAR}.

\textbf{Lote en sem. 2:} cubre semanas 2 y 3.
\(Q_2 = 30 + 120 = \mathbf{150}\).

\textbf{Lote 2 --- arranca en período 4 (\(D_4 = 180\)):} - \(k=1\):
\(C(1) = 120\). - \(k=2\) (sem. 4--5, \(D_5 = 150\)):
\(C(2) = (120 + 1 \cdot 0{,}9 \cdot 150)/2 = 127{,}5 > 120\) →
\textbf{PARAR}.

\textbf{Lote en sem. 4:} cubre solo semana 4. \(Q_4 = \mathbf{180}\).

\textbf{Lote 3 --- arranca en período 5 (\(D_5 = 150\)):} - Único
período restante con demanda. \textbf{Lote en sem. 5:}
\(Q_5 = \mathbf{150}\).

\textbf{Política resultante:}

\begin{longtable}[]{@{}lcccccccc@{}}
\toprule\noalign{}
Período & 1 & 2 & 3 & 4 & 5 & 6 & 7 & 8 \\
\midrule\noalign{}
\endhead
\bottomrule\noalign{}
\endlastfoot
\textbf{Demanda} & 0 & 30 & 120 & 180 & 150 & 0 & 0 & 0 \\
\textbf{Lote \(Q\)} & 0 & \textbf{150} & 0 & \textbf{180} & \textbf{150}
& 0 & 0 & 0 \\
\textbf{Inventario} & 0 & 120 & 0 & 0 & 0 & 0 & 0 & 0 \\
\end{longtable}

\textbf{Costos totales:} - Setups: \(3 \times \$120 = \$360\). -
Almacenamiento: 120 unidades almacenadas 1 semana en semana 3 →
\(120 \times 1 \times 0{,}9 = \$108\).

\[\boxed{\text{Costo Total Silver-Meal} = 360 + 108 = \$468}\]

\begin{center}\rule{0.5\linewidth}{0.5pt}\end{center}

\hypertarget{problema-3-pert-y-crashing}{%
\subsection{Problema 3: PERT y
Crashing}\label{problema-3-pert-y-crashing}}

\hypertarget{enunciado-2}{%
\subsubsection{Enunciado}\label{enunciado-2}}

\begin{longtable}[]{@{}ccccc@{}}
\toprule\noalign{}
Etapa & Predecesor & \(a\) (optim.) & \(m\) (probable) & \(b\)
(pesim.) \\
\midrule\noalign{}
\endhead
\bottomrule\noalign{}
\endlastfoot
A & --- & 2 & 3 & 4 \\
B & A & 2 & 4 & 6 \\
C & A & 5 & 6 & 13 \\
D & B, C & 3 & 6 & 9 \\
E & B & 2 & 5 & 8 \\
F & D, E & 2 & 4 & 6 \\
\end{longtable}

\textbf{Datos de crashing} (costo acelerado = asumir antes de comenzar,
varianza inalterada):

\begin{longtable}[]{@{}
  >{\centering\arraybackslash}p{(\columnwidth - 8\tabcolsep) * \real{0.2000}}
  >{\centering\arraybackslash}p{(\columnwidth - 8\tabcolsep) * \real{0.2000}}
  >{\centering\arraybackslash}p{(\columnwidth - 8\tabcolsep) * \real{0.2000}}
  >{\centering\arraybackslash}p{(\columnwidth - 8\tabcolsep) * \real{0.2000}}
  >{\centering\arraybackslash}p{(\columnwidth - 8\tabcolsep) * \real{0.2000}}@{}}
\toprule\noalign{}
\begin{minipage}[b]{\linewidth}\centering
Actividad
\end{minipage} & \begin{minipage}[b]{\linewidth}\centering
Reducción máx. (sem.)
\end{minipage} & \begin{minipage}[b]{\linewidth}\centering
Costo Normal (\$)
\end{minipage} & \begin{minipage}[b]{\linewidth}\centering
Costo Acelerado (\$)
\end{minipage} & \begin{minipage}[b]{\linewidth}\centering
Δ Costo
\end{minipage} \\
\midrule\noalign{}
\endhead
\bottomrule\noalign{}
\endlastfoot
A & 1 & 10 000 & 13 000 & 3 000 \\
B & 1 & 6 000 & 9 000 & 3 000 \\
C & 2 & 4 000 & 7 000 & 3 000 \\
D & 2 & 13 000 & 18 000 & 5 000 \\
E & 2 & 9 000 & 13 000 & 4 000 \\
F & 1 & 7 000 & 8 000 & 1 000 \\
\end{longtable}

\textbf{Se pide:} - (a) Tiempos esperados y varianzas. - (b) Diagrama
PERT. - (c) ES, EF, LS, LF, holgura, ruta crítica y duración mínima. -
(d) Probabilidad de terminar en \textless{} 22 semanas. IC 95\%. - (e)
Si el cliente ofrece bono de \(\$8\,000\) al terminar en \textless{} 18
semanas, ¿qué actividades acortaría?

\hypertarget{soluciuxf3n-oficial-2}{%
\subsubsection{Solución Oficial}\label{soluciuxf3n-oficial-2}}

\hypertarget{a-tiempos-esperados-y-varianzas}{%
\paragraph{(a) Tiempos Esperados y
Varianzas}\label{a-tiempos-esperados-y-varianzas}}

\[t_e = \frac{a + 4m + b}{6}, \qquad \sigma^2 = \left(\frac{b-a}{6}\right)^2\]

\begin{longtable}[]{@{}ccc@{}}
\toprule\noalign{}
Actividad & \(t_e\) (sem.) & \(\sigma^2\) \\
\midrule\noalign{}
\endhead
\bottomrule\noalign{}
\endlastfoot
A & 3 & 0,111 (≈ 1/9) \\
B & 4 & 0,444 (≈ 4/9) \\
C & 7 & 1,778 (≈ 16/9) \\
D & 6 & 1,000 \\
E & 5 & 1,000 \\
F & 4 & 0,444 (≈ 4/9) \\
\end{longtable}

\begin{quote}
Verificación C: \(t_e = (5 + 4 \times 6 + 13)/6 = 42/6 = 7\).
\(\sigma^2 = ((13-5)/6)^2 = (8/6)^2 = 64/36 = 1{,}78\).
\end{quote}

\begin{center}\rule{0.5\linewidth}{0.5pt}\end{center}

\hypertarget{b-diagrama-pert}{%
\paragraph{(b) Diagrama PERT}\label{b-diagrama-pert}}

\begin{verbatim}
         B(4)           E(5)
    2 ────────── 3 ────────── 5
   /                           \
1 ─── A(3)                       ─── F(4) ─── 6
   \                           /
    ─────────── 4 ────────── 5
         C(7)           D(6)
\end{verbatim}

Nodos: 1 (inicio) → {[}A{]} → 2 → {[}B{]} → 3 → {[}E{]} → 5 → {[}F{]} →
6\\
1 → {[}A{]} → 2 → {[}C{]} → 4 → {[}D{]} → 5 → {[}F{]} → 6

\begin{center}\rule{0.5\linewidth}{0.5pt}\end{center}

\hypertarget{c-tiempos-de-calendario-y-ruta-cruxedtica}{%
\paragraph{(c) Tiempos de Calendario y Ruta
Crítica}\label{c-tiempos-de-calendario-y-ruta-cruxedtica}}

\textbf{Pasada forward:}

\begin{longtable}[]{@{}cccc@{}}
\toprule\noalign{}
Actividad & \(t_e\) & ES & EF \\
\midrule\noalign{}
\endhead
\bottomrule\noalign{}
\endlastfoot
A & 3 & 0 & 3 \\
B & 4 & 3 & 7 \\
C & 7 & 3 & 10 \\
D & 6 & max(7, 10) = \textbf{10} & 16 \\
E & 5 & 7 & 12 \\
F & 4 & max(16, 12) = \textbf{16} & \textbf{20} \\
\end{longtable}

\textbf{Duración del proyecto: T = 20 semanas.}

\textbf{Pasada backward} (\(LF_F = 20\)):

\begin{longtable}[]{@{}ccccc@{}}
\toprule\noalign{}
Actividad & \(t_e\) & LF & LS & \textbf{Holgura} \\
\midrule\noalign{}
\endhead
\bottomrule\noalign{}
\endlastfoot
F & 4 & 20 & 16 & \textbf{0} \\
E & 5 & 16 & 11 & \textbf{4} \\
D & 6 & 16 & 10 & \textbf{0} \\
C & 7 & 10 & 3 & \textbf{0} \\
B & 4 & min(10, 11) = 10 & 6 & \textbf{3} \\
A & 3 & min(6, 3) = 3 & 0 & \textbf{0} \\
\end{longtable}

\textbf{Ruta crítica: A → C → D → F.} Holgura = 0 en todas.

\begin{center}\rule{0.5\linewidth}{0.5pt}\end{center}

\hypertarget{d-probabilidades-e-intervalo-de-confianza}{%
\paragraph{(d) Probabilidades e Intervalo de
Confianza}\label{d-probabilidades-e-intervalo-de-confianza}}

\[\sigma_p = \sqrt{\sigma_A^2 + \sigma_C^2 + \sigma_D^2 + \sigma_F^2} = \sqrt{0{,}111 + 1{,}778 + 1{,}000 + 0{,}444} = \sqrt{3{,}333} \approx 1{,}82 \text{ sem.}\]

\textbf{Probabilidad de terminar en \textless{} 22 semanas:}
\[Z = \frac{22 - 20}{1{,}82} \approx 1{,}10 \implies \Phi(1{,}10) = 0{,}864 \implies \mathbf{86{,}4\%}\]

\textbf{Intervalo de confianza al 95\% (bilateral, \(Z = 1{,}96\)):}
\[IC_{95\%} = 20 \pm 1{,}96 \times 1{,}82 = 20 \pm 3{,}57 \implies [16{,}43;\; 23{,}57] \text{ semanas}\]

\begin{center}\rule{0.5\linewidth}{0.5pt}\end{center}

\hypertarget{e-crashing-para-bono-de-8000-18-semanas}{%
\paragraph{\texorpdfstring{(e) Crashing para Bono de \(\$8\,000\)
(\textless{} 18
semanas)}{(e) Crashing para Bono de \textbackslash\$8\textbackslash,000 (\textless{} 18 semanas)}}\label{e-crashing-para-bono-de-8000-18-semanas}}

Necesitamos reducir \textbf{2 semanas} en la ruta crítica A--C--D--F. Se
evalúa el \textbf{valor esperado neto} para cada escenario.

\textbf{Escenario 0: Sin aceleración (duración = 20 sem.)}
\[Z = \frac{18 - 20}{1{,}82} = -1{,}10 \implies P(T < 18) = 13{,}6\%\]
\[VE_{\text{neto}} = 0{,}136 \times 8\,000 - 0 = \$1\,086\]

\textbf{Escenario 1: Acortar 1 semana --- opciones en ruta crítica:}

\begin{longtable}[]{@{}ccc@{}}
\toprule\noalign{}
Actividad & Reducción & Costo \\
\midrule\noalign{}
\endhead
\bottomrule\noalign{}
\endlastfoot
A & 1 sem. & \(\$3\,000\) \\
F & 1 sem. & \(\$1\,000\) \\
\end{longtable}

Óptimo: acelerar \textbf{F} (\(\$1\,000\)). Nueva duración = 19 sem.
\[Z = \frac{18 - 19}{1{,}82} = -0{,}55 \implies P(T < 18) = 29{,}1\%\]
\[VE_{\text{neto}} = 0{,}291 \times 8\,000 - 1\,000 = 2\,328 - 1\,000 = \$1\,330\]

\textbf{Escenario 2: Acortar 2 semanas --- opciones en ruta crítica:}

\begin{longtable}[]{@{}lc@{}}
\toprule\noalign{}
Combinación & Costo total \\
\midrule\noalign{}
\endhead
\bottomrule\noalign{}
\endlastfoot
A (1 sem.) + F (1 sem.) & \(\$4\,000\) \\
\textbf{C (2 sem.)} & \textbf{\(\$3\,000\)} \\
D (2 sem.) & \(\$5\,000\) \\
\end{longtable}

Óptimo: acelerar \textbf{C} en 2 semanas (\(\$3\,000\)). Nueva duración
= 18 sem. \[Z = \frac{18 - 18}{1{,}82} = 0 \implies P(T < 18) = 50\%\]
\[VE_{\text{neto}} = 0{,}50 \times 8\,000 - 3\,000 = 4\,000 - 3\,000 = \$1\,000\]

\textbf{Resumen comparativo:}

\begin{longtable}[]{@{}
  >{\raggedright\arraybackslash}p{(\columnwidth - 8\tabcolsep) * \real{0.1667}}
  >{\centering\arraybackslash}p{(\columnwidth - 8\tabcolsep) * \real{0.2083}}
  >{\centering\arraybackslash}p{(\columnwidth - 8\tabcolsep) * \real{0.2083}}
  >{\centering\arraybackslash}p{(\columnwidth - 8\tabcolsep) * \real{0.2083}}
  >{\centering\arraybackslash}p{(\columnwidth - 8\tabcolsep) * \real{0.2083}}@{}}
\toprule\noalign{}
\begin{minipage}[b]{\linewidth}\raggedright
Escenario
\end{minipage} & \begin{minipage}[b]{\linewidth}\centering
Costo aceleración
\end{minipage} & \begin{minipage}[b]{\linewidth}\centering
\(P(T < 18)\)
\end{minipage} & \begin{minipage}[b]{\linewidth}\centering
VE bruto
\end{minipage} & \begin{minipage}[b]{\linewidth}\centering
\textbf{VE neto}
\end{minipage} \\
\midrule\noalign{}
\endhead
\bottomrule\noalign{}
\endlastfoot
Sin aceleración & \(\$0\) & 13,6\% & \(\$1\,086\) & \(\$1\,086\) \\
Acelerar F 1 sem. & \(\$1\,000\) & 29,1\% & \(\$2\,328\) &
\textbf{\(\$1\,330\)} ← máximo \\
Acelerar C 2 sem. & \(\$3\,000\) & 50,0\% & \(\$4\,000\) &
\(\$1\,000\) \\
\end{longtable}

\[\boxed{\text{Decisión: Acelerar F 1 semana (costo } \$1\,000\text{). VE neto máximo = } \$1\,330}\]

\begin{quote}
\textbf{Verificación: ¿se vuelve crítica alguna ruta no crítica tras el
crashing?} Las rutas alternativas y sus duraciones originales son: -
A--B--D--F: \(3+4+6+4 = 17\) sem. - A--B--E--F: \(3+4+5+4 = 16\) sem.

Tras acelerar F en 1 semana, la nueva duración de F = 3 sem. Las rutas
alternativas pasan a: - A--B--D--F: \(3+4+6+3 = 16\) sem. \(< 19\) ✓ -
A--B--E--F: \(3+4+5+3 = 15\) sem. \(< 19\) ✓

Ninguna supera la nueva RC (19 sem.), por lo que el análisis es válido y
no se generan rutas críticas adicionales.
\end{quote}

\begin{center}\rule{0.5\linewidth}{0.5pt}\end{center}

\hypertarget{problema-4-variabilidad-optimizaciuxf3n-de-capacidad}{%
\subsection{Problema 4: Variabilidad --- Optimización de
Capacidad}\label{problema-4-variabilidad-optimizaciuxf3n-de-capacidad}}

\hypertarget{enunciado-3}{%
\subsubsection{Enunciado}\label{enunciado-3}}

Empresa de imprenta: clientes llegan a tasa \(\lambda\)
{[}clientes/hr{]}, capacidad de atención \(\mu\) {[}clientes/hr{]},
distribuciones generales con tiempo efectivo \(t_e\), coeficientes de
variación \(c_a\) (llegadas) y \(c_e\) (servicio).

La empresa evalúa ofrecer descuentos \(\Delta\) para modular la demanda:
\[\lambda(\Delta) = \lambda_0 e^{-\Delta}\] El costo de espera del
cliente es \(CE(W) = 100 + W\), donde \(W\) es el tiempo total en el
sistema.

\textbf{a)} Plantear el modelo de programación matemática y las
condiciones de primer orden.

\textbf{b)} Si se puede adquirir capacidad adicional \(IC\)
{[}clientes/hr{]} a un costo \(\$K\) por unidad, ¿cómo cambia el modelo?

\hypertarget{soluciuxf3n-oficial-3}{%
\subsubsection{Solución Oficial}\label{soluciuxf3n-oficial-3}}

\hypertarget{a-modelo-de-optimizaciuxf3n-variable-de-decisiuxf3n-delta}{%
\paragraph{\texorpdfstring{a) Modelo de Optimización (variable de
decisión:
\(\Delta\))}{a) Modelo de Optimización (variable de decisión: \textbackslash Delta)}}\label{a-modelo-de-optimizaciuxf3n-variable-de-decisiuxf3n-delta}}

\[\min_{\Delta \ge 0} \quad F = 100 + W + \Delta\]

Sujeto a:
\[W = \left(\frac{c_a^2 + c_e^2}{2}\right) \cdot \frac{\rho}{1 - \rho} \cdot t_e, \qquad \rho = \frac{\lambda}{\mu} = \frac{\lambda_0 e^{-\Delta}}{\mu}, \qquad \Delta \ge 0\]

\textbf{Forma compacta} (sustituyendo
\(\lambda = \lambda_0 e^{-\Delta}\) directamente):

\[\min_{\Delta \ge 0} \quad F = 100 + \left(\frac{c_a^2 + c_e^2}{2}\right) \cdot \frac{\dfrac{\lambda_0 e^{-\Delta}}{\mu}}{1 - \dfrac{\lambda_0 e^{-\Delta}}{\mu}} \cdot t_e \;+\; \Delta\]

\textbf{Condición de primer orden}
(\(\partial F / \partial \Delta = 0\), no resolver):
\[\frac{\partial F}{\partial \Delta} = \frac{\partial W}{\partial \Delta} + 1 = 0\]

donde \(\partial W / \partial \Delta\) involucra la derivada de
\(\rho/(1-\rho)\) respecto a \(\Delta\), que resulta negativa (al
aumentar \(\Delta\), baja \(\lambda\) y baja la espera). La CPO iguala
la ganancia marginal en espera con el costo marginal del descuento.

\begin{center}\rule{0.5\linewidth}{0.5pt}\end{center}

\hypertarget{b-modelo-con-capacidad-adicional-variables-de-decisiuxf3n-delta-y-ic}{%
\paragraph{\texorpdfstring{b) Modelo con Capacidad Adicional (variables
de decisión: \(\Delta\) y
\(IC\))}{b) Modelo con Capacidad Adicional (variables de decisión: \textbackslash Delta y IC)}}\label{b-modelo-con-capacidad-adicional-variables-de-decisiuxf3n-delta-y-ic}}

Se suma \(IC\) a la capacidad base, por lo que
\(\rho = \lambda / (\mu + IC)\):

\[\min_{\Delta \ge 0,\; IC \ge 0} \quad F = 100 + W(\Delta, IC) + \Delta + K \cdot IC\]

Sujeto a:
\[W = \left(\frac{c_a^2 + c_e^2}{2}\right) \cdot \frac{\rho}{1 - \rho} \cdot t_e, \qquad \rho = \frac{\lambda_0 e^{-\Delta}}{\mu + IC}\]
\[\lambda_0 e^{-\Delta} < \mu + IC \quad (\text{condición de estabilidad})\]
\[\Delta \ge 0, \quad IC \ge 0\]

\begin{quote}
\textbf{Interpretación:} Ahora hay dos palancas para reducir la espera:
bajar la demanda vía descuento (\(\Delta\)) o aumentar la capacidad
(\(IC\)). La CPO para \(IC\) iguala el costo marginal \(K\) con la
reducción marginal en \(W\).
\end{quote}

\end{document}
